\documentclass[../../main.tex]{subfiles}% 注意这里的文件路径不能用 ./main.tex ,否则用latexmk编译子文件会报错
\graphicspath{{\subfix{./image/}}{\subfix{../../image/}}} % 指定图片引用路径,后续可以直接使用图片文件名
% 如果图片存放在上级目录,则使用 ../../image/ 作为备用路径

% 例如:
% \begin{figure}[H]
% \centering
% \includegraphics[width=0.3\textwidth]{图.png}
% \caption{}
% \label{figure:图}
% \end{figure}
% 注意:上述\label{}一定要放在\caption{}之后,否则引用图片序号会只会显示??.

\begin{document}

\section{基本解和 Green 函数}

\subsection{基本解}

\begin{definition}
对 $x\in\mathbb{R}^n$，$x\neq 0$，称函数
\begin{align}
\varGamma (x) = \begin{cases}
-\frac{1}{2\pi}\ln|x|, & n=2, \\
\frac{1}{n(n-2)\alpha(n)|x|^{n-2}}, & n\geqslant 3
\end{cases} \label{eq:fundamental}
\end{align}
为\textbf{Laplace 方程的基本解},其中 $\alpha(n)$ 表示 $\mathbb{R}^n$ 中单位球的体积.
\end{definition}
\begin{remark}
Laplace方程的基本解具有很强的物理意义：当 $n=3$ 时，基本解 $\varGamma (x)$ 就是由放置在原点的单位正点电荷引起的在全空间 $\mathbb{R}^3$ 上的静电位势分布。
\end{remark}

\begin{proposition}\label{proposition:Laplace方程基本解的性质}
设 $\varGamma (x)$ 是 $\mathbb{R}^n$ 上 Laplace 方程的基本解，则 $\varGamma $ 具有以下性质：
\begin{enumerate}[(1)]
\item\label{proposition:Laplace方程基本解的性质-1} $\varGamma \in L^1_{\mathrm{loc}}(\mathbb{R}^n)$，即 $\varGamma $ 在 $\mathbb{R}^n$ 上局部Lebesgue可积；

\item\label{proposition:Laplace方程基本解的性质-2} 对任意 $x\neq 0$，存在仅依赖于 $n$ 的常数 $C>0$，使得
\begin{align*}
|D\varGamma (x)| \leqslant \frac{C}{|x|^{n-1}},\qquad |D^2\varGamma (x)| \leqslant \frac{C}{|x|^{n}}.
\end{align*}

\item\label{proposition:Laplace方程基本解的性质-3} 在广义函数意义下，$\varGamma $ 满足
\begin{align*}
-\Delta \varGamma  = \delta,
\end{align*}
其中 $\delta$ 是原点处的 Dirac 测度. 等价地，对任意测试函数 $f\in C_0^\infty(\mathbb{R}^n)$，有
\begin{align*}
\int_{\mathbb{R}^n} \varGamma (x)(-\Delta f(x))\,\mathrm{d}x = f(0). 
\end{align*}
\end{enumerate}
\end{proposition}
\begin{remark}
性质\ref{proposition:Laplace方程基本解的性质-2}中的常数 $C$ 可能在不同估计式中取不同值，但始终仅依赖于 $n$.
\end{remark}

\begin{theorem}\label{theorem:Laplace方程基本解的性质}
假设 $f(x)\in C_0^2(\mathbb{R}^n)$，定义
\begin{align*}
u(x) = \int_{\mathbb{R}^n} \varGamma (x-y)f(y)\,\mathrm{d}y,
\end{align*}
其中$\varGamma (x)$是Laplace方程的基本解.则 $u\in C^2(\mathbb{R}^n)$ 且满足 $-\Delta u=f$。
\end{theorem}
\begin{remark}
此定理说明我们可以利用基本解来构造位势方程(Poisson 方程)在全空间上的解。
\end{remark}
\begin{proof}
固定 $x\in\mathbb{R}^n$。
我们注意到
\begin{align*}
u(x) = \int_{\mathbb{R}^n} \varGamma (y)f(x-y)\,\mathrm{d}y,
\end{align*}
从而有
\begin{align*}
\frac{u(x+he_i)-u(x)}{h} = \int_{\mathbb{R}^n} \varGamma (y)\left[ \frac{f(x+he_i-y)-f(x-y)}{h} \right] \mathrm{d}y,
\end{align*}
其中 $e_i$ 表示 $\mathbb{R}^n$ 上的第 $i$ 个方向向量。由于 $f$ 具有紧支集，则
\begin{align*}
\lim_{h\to 0} \frac{f(x+he_i-y)-f(x-y)}{h} = \frac{\partial}{\partial x_i} f(x-y),
\end{align*}
且上述极限关于 $y$ 一致收敛。于是
\begin{align*}
\frac{\partial u}{\partial x_i} = \int_{\mathbb{R}^n} \varGamma (y) \frac{\partial}{\partial x_i} f(x-y)\,\mathrm{d}y \quad (i=1,2,\cdots,n).
\end{align*}
同理我们得到
\begin{align*}
\frac{\partial^2 u}{\partial x_i\partial x_j} = \int_{\mathbb{R}^n} \varGamma (y) \frac{\partial^2}{\partial x_i\partial x_j} f(x-y)\,\mathrm{d}y \quad (i,j=1,2,\cdots,n).
\end{align*}
由于上式右端项对变量 $x$ 是连续的，因而 $u\in C^2(\mathbb{R}^n)$。实际上，由于 $f$ 具有紧支集，我们有：存在 $M_0>0$，使得对所有 $x\in\mathbb{R}^n$，有
\begin{align*}
|f(x)| + |Df(x)| + |D^2f(x)| \leqslant M_0.
\end{align*}
由此容易证明 $u, Du, D^2u$ 是有界的，即存在 $M_1>0$ 使得对所有 $x\in\mathbb{R}^n$，有
\begin{align*}
|u(x)| + |Du(x)| + |D^2u(x)| \leqslant M_1.
\end{align*}

固定 $\varepsilon>0$，则
\begin{align*}
\Delta u &= \int_{\mathbb{R}^n} \varGamma (y)\Delta_x f(x-y)\,\mathrm{d}y = \int_{B(0,\varepsilon)} \varGamma (y)\Delta_x f(x-y)\,\mathrm{d}y + \int_{\mathbb{R}^n\setminus B(0,\varepsilon)} \varGamma (y)\Delta_x f(x-y)\,\mathrm{d}y,
\end{align*}
这里 $\Delta_x$ 是关于变量 $x$ 的 Laplace 算子 $\Delta$。记
\begin{align*}
I_\varepsilon = \int_{B(0,\varepsilon)} \varGamma (y)\Delta_x f(x-y)\,\mathrm{d}y, \quad J_\varepsilon = \int_{\mathbb{R}^n\setminus B(0,\varepsilon)} \varGamma (y)\Delta_x f(x-y)\,\mathrm{d}y,
\end{align*}
于是
\begin{align*}
\Delta u = I_\varepsilon + J_\varepsilon.
\end{align*}

我们首先估计 $I_\varepsilon$。利用函数 $f$ 的有界性我们容易证明
\begin{align}
|I_\varepsilon| \leqslant M_0 \int_{B(0,\varepsilon)} |\varGamma (y)|\,\mathrm{d}y
\leqslant \begin{cases}
C\varepsilon^2|\ln\varepsilon|, & n=2, \\
C\varepsilon^2, & n\geqslant 3,
\end{cases} \label{eq:Ieps}
\end{align}
其中 $C$ 是不依赖于 $\varepsilon$ 的常数。

由于函数 $f(y)\in C_0^2(\mathbb{R}^n)$ 具有紧支集，因而函数 $f(x-y)$ 具有紧支集。故存在一个正数 $R_x>0$，使得 $f(x-y)$ 的支集 $\operatorname{spt} f(x-y)\subset B(0,R_x)$。由\hyperref[Basis of Analytics-theorem:Green第一恒等式]{Green第一恒等式} 我们可以将 $J_\varepsilon$ 表示成
\begin{align*}
J_\varepsilon &= \int_{\mathbb{R}^n\setminus B(0,\varepsilon)} \varGamma (y)\Delta_y f(x-y)\,\mathrm{d}y 
= \int_{B(0,R_x)\setminus B(0,\varepsilon)} \varGamma (y)\Delta_y f(x-y)\,\mathrm{d}y \\
&= \int_{B(0,R_x)\setminus B(0,\varepsilon)} \Delta\varGamma (y) f(x-y)\,\mathrm{d}y + \int_{\partial B(0,R_x)\cup\partial B(0,\varepsilon)} \left[ \varGamma (y)\frac{\partial}{\partial n} f(x-y) - \frac{\partial\varGamma (y)}{\partial n} f(x-y) \right] \mathrm{d}S(y) \\
&= \int_{\partial B(0,\varepsilon)} \left[ \varGamma (y)\frac{\partial}{\partial n} f(x-y) - \frac{\partial\varGamma (y)}{\partial n} f(x-y) \right] \mathrm{d}S(y),
\end{align*}
其中在球面 $\partial B(0,\varepsilon)$ 上 $n$ 表示单位内法向量，在球面 $\partial B(0,R_x)$ 上 $n$ 表示单位外法向量。记
\begin{align*}
K_\varepsilon = \int_{\partial B(0,\varepsilon)} \varGamma (y)\frac{\partial}{\partial n} f(x-y)\,\mathrm{d}S(y), \quad 
L_\varepsilon = -\int_{\partial B(0,\varepsilon)} \frac{\partial\varGamma (y)}{\partial n} f(x-y)\,\mathrm{d}S(y),
\end{align*}
于是
\begin{align*}
J_\varepsilon = K_\varepsilon + L_\varepsilon.
\end{align*}

我们首先估计 $K_\varepsilon$，然后再计算 $L_\varepsilon$。利用函数 $f$ 的梯度 $Df$ 的有界性我们容易证明
\begin{align}
|K_\varepsilon| \leqslant M_0 \int_{\partial B(0,\varepsilon)} |\varGamma (y)|\,\mathrm{d}S(y)
\leqslant \begin{cases}
C\varepsilon|\ln\varepsilon|, & n=2, \\
C\varepsilon, & n\geqslant 3,
\end{cases} \label{eq:Keps}
\end{align}
其中 $C$ 是不依赖于 $\varepsilon$ 的常数。

注意到在 $\partial B(0,\varepsilon)$ 上，有
\begin{align*}
D\varGamma (y) = \frac{-y}{n\alpha(n)|y|^n}, \qquad n = -\frac{y}{|y|} = -\frac{y}{\varepsilon},
\end{align*}
于是
\begin{align*}
\frac{\partial\varGamma (y)}{\partial n} = D\varGamma (y)\cdot n = \frac{1}{n\alpha(n)\varepsilon^{n-1}}.
\end{align*}
由中值公式我们得到
\begin{align}
L_\varepsilon = -\frac{1}{n\alpha(n)\varepsilon^{n-1}} \int_{\partial B(0,\varepsilon)} f(x-y)\,\mathrm{d}S(y)
= -\frac{f(x-y_\varepsilon)}{n\alpha(n)\varepsilon^{n-1}} \int_{\partial B(0,\varepsilon)} \mathrm{d}S(y) = -f(x-y_\varepsilon), \label{eq:Leps}
\end{align}
其中 $y_\varepsilon\in\partial B(0,\varepsilon)$。由函数 $f$ 的连续性，当 $\varepsilon\to 0$ 时，$L_\varepsilon\to -f(x)$。

综上我们知道，对任意 $\varepsilon>0$，有
\begin{align*}
\Delta u = I_\varepsilon + K_\varepsilon + L_\varepsilon.
\end{align*}
利用估计 \eqref{eq:Ieps}, \eqref{eq:Keps} 和 \eqref{eq:Leps}，令 $\varepsilon\to 0$，则有
\begin{align*}
-\Delta u(x) = f(x).
\end{align*}
至此我们完成定理的证明.
\end{proof}

\begin{theorem}
假设 $f(x)\in C_0^2(\mathbb{R}^n)$，则
\begin{align*}
u(x) = \int_{\mathbb{R}^n} \varGamma (x-y)f(y)\,\mathrm{d}y + C
\end{align*}
是位势方程在全空间 $\mathbb{R}^n$ 上所有的有界解。
\end{theorem}
\begin{proof}
利用\hyperref[theorem:Laplace方程基本解的性质]{Laplace方程基本解的性质}和 \hyperref[theorem:PDE-Liouville定理]{Liouville定理}得证.
\end{proof}

\begin{definition}
设$\Omega \subset \mathbb{R}^n$,齐次边值问题
\begin{align}
\begin{cases}
-\Delta u = \lambda u, & x\in\Omega, \\
\left.\frac{\partial u}{\partial n}\right|_{\partial\Omega} = 0
\end{cases} \label{eq:eigenvalue}
\end{align}
称为\textbf{特征值问题}，使此问题有非零解的 $\lambda\in\mathbb{R}$ 称为此问题的\textbf{特征值}，相应的非零解称为对应于这个特征值的\textbf{特征函数}，记为 $u_\lambda$。
\end{definition}
\begin{remark}
在特征值问题 \eqref{eq:eigenvalue} 的方程两端同乘 $u$ 并在 $\Omega$ 上积分，我们不难证明特征值 $\lambda\geqslant 0$。
\end{remark}

\begin{theorem}
设$\Omega \subset \mathbb{R}^n$,则第二边值问题
\begin{align}
\begin{cases}
-\Delta u = \lambda u + f, & x\in\Omega, \\
\left.\frac{\partial u}{\partial n}\right|_{\partial\Omega} = g
\end{cases} \label{eq:neumann-gen}
\end{align}
有 $C^1(\overline{\Omega})\cap C^2(\Omega)$ 上的解的必要条件是对于相应的任意特征函数 $u_\lambda(x)$，成立
\begin{align}
\int_\Omega f(x)u_\lambda(x)\,\mathrm{d}x + \int_{\partial\Omega} g(x)u_\lambda(x)\,\mathrm{d}S(x) = 0. \label{eq:neumann-cond}
\end{align}
\end{theorem}

\begin{proof}
在\hyperref[Basis of Analytics-theorem:Green第一恒等式]{Green第一恒等式} 公式中，取 $u$ 为第二边值问题 \eqref{eq:neumann-gen} 的解，$v=u_\lambda(x)$，立即得到要证的结论.
\end{proof}

\begin{corollary}
设$\Omega \subset \mathbb{R}^n$,则Poisson 方程的 Neumann 边值问题
\begin{align}
\begin{cases}
-\Delta u = f, & x\in\Omega, \\
\left.\frac{\partial u}{\partial n}\right|_{\partial\Omega} = g
\end{cases} \label{eq:neumann}
\end{align}
有 $C^1(\overline{\Omega})\cap C^2(\Omega)$ 上的解的必要条件是
\begin{align}
\int_\Omega f(x)\,\mathrm{d}x + \int_{\partial\Omega} g(x)\,\mathrm{d}S(x) = 0. \label{eq:neumann-cond0}
\end{align}
\end{corollary}
\begin{proof}
实际上，对于 Neumann 边值问题 \eqref{eq:neumann} 的方程两端在 $\Omega$ 上积分，利用\hyperref[Basis of Analytics-theorem:Gauss-Green公式(散度定理)]{Gauss-Green公式}也可以得到上式.
\end{proof}


\subsection{Green 函数}

\begin{definition}
设$\Omega \subset \mathbb{R}^n$,对给定$x\in \Omega$,函数$\phi^x$满足方程
\begin{align*}
\begin{cases}
-\Delta \phi ^x\left( y \right) =0,\quad y\in \Omega ,\\
\phi ^x\mid _{\partial \Omega}=\varGamma  \left( y-x \right) .\\
\end{cases}
\end{align*}
其中$\varGamma (x)$是Laplace方程的基本解.对于任意 $x,y\in\Omega$，$x\neq y$，函数
\begin{align*}
G(x,y) = \varGamma (y-x) - \phi^x(y)
\end{align*}
称为 $\Omega$ 上的 \textbf{Green 函数}.
\end{definition}
\begin{remark}
显然，当 $x\in\Omega$，$y\in\partial\Omega$ 时，$G(x,y)=0$。实际上，Green 函数就是基本解减去一个以基本解为边值的调和函数。由于函数 $\varGamma (y-x)$ 在 $\partial\Omega$ 附近是一个 $C^\infty$ 的光滑函数，当 $\Omega$ 满足一定的条件时，$\phi^x(y)$ 的存在性可以比较容易地得到。
\end{remark}
\begin{remark}
Green 函数 $G(x,y)$ $(n=2,3)$ 的物理意义如下：在物体内部 $x$ 处放置一个单位点热源，与外界接触的表面保持零度，那么物体的稳定温度场就是 Green 函数；或者让某导体的表面接地，在其内部 $x$ 处放置一个单位正电荷，那么在导体内部所产生的电势分布也是 Green 函数。
\end{remark}
\begin{remark}
引入 Green 函数 $G(x,y)$ 的重要意义在于把求解具有任意非齐次项与任意边值的定解问题  归结为求解一个特定的边值问题。对于一些特殊区域，这样的特定的边值问题可以得到解的具体表达式。在一般情形，虽然不可能给出 Green 函数的表达式，但 Green 函数只依赖于区域，而与边值和非齐次项无关，这无论对理论研究还是对求解问题都带来很大的方便。
\end{remark}

\begin{theorem}\label{theorem:Dirichlet问题解的表达式}
假设 $\Omega$ 是 $\mathbb{R}^n$ 上的一个有界区域，$u(x)\in C^2(\Omega)\cap C^1(\overline{\Omega})$ 是 Dirichlet 问题
\begin{align*}
\begin{cases}
-\Delta u = f, & x\in\Omega, \\
u = g, & x\in\partial\Omega
\end{cases} 
\end{align*}
的解，则
\begin{align}
u(x) = -\int_{\partial\Omega} g(y)\frac{\partial}{\partial n}G(x,y)\,\mathrm{d}S(y) + \int_\Omega G(x,y)f(y)\,\mathrm{d}y, \label{eq:green-sol}
\end{align}
其中$G(x,y)$是$\Omega$上的Green函数.
\end{theorem}
\begin{proof}
固定 $x\in\Omega$. 取 $\varepsilon>0$ 充分小使得 $B(x,\varepsilon)\subset\Omega$. 在区域 $\Omega_\varepsilon\triangleq\Omega\setminus \overline{B(x,\varepsilon)}$ 上对 $u(y)$ 和 $\varGamma (y-x)$ 应用\hyperref[Basis of Analytics-theorem:Green第二恒等式]{Green第二恒等式}，得
\begin{align}
\int_{\Omega_\varepsilon} \big[u(y)\Delta\varGamma (y-x)-\varGamma (y-x)\Delta u(y)\big]\,\mathrm{d}y
= \int_{\partial\Omega} \left[u(y)\frac{\partial}{\partial n}\varGamma (y-x)-\varGamma (y-x)\frac{\partial u}{\partial n}(y)\right]\mathrm{d}S(y) \nonumber\\
\quad + \int_{\partial B(x,\varepsilon)} \left[u(y)\frac{\partial}{\partial n}\varGamma (y-x)-\varGamma (y-x)\frac{\partial u}{\partial n}(y)\right]\mathrm{d}S(y), \label{eq:green2}
\end{align}
其中在 $\partial\Omega$ 上 $n$ 为$\partial\Omega$ 上的单位外法向量，在 $\partial B(x,\varepsilon)$ 上 $n$ 为$B(x,\varepsilon)$的单位外法向量.

由于在 $\Omega_\varepsilon$ 上 $\Delta\varGamma (y-x)=0$，且 $\Delta u=-f$，故 \eqref{eq:green2}式左边等于
\begin{align*}
\int_{\Omega_\varepsilon} \varGamma (y-x) f(y)\,\mathrm{d}y.
\end{align*}
令 $\varepsilon\to 0$，有
\begin{align*}
\int_{\partial B(x,\varepsilon)} u(y)\frac{\partial}{\partial n}\varGamma (y-x)\,\mathrm{d}S(y) \longrightarrow u(x),
\end{align*}
以及
\begin{align*}
\int_{\partial B(x,\varepsilon)} \varGamma (y-x)\frac{\partial u}{\partial n}(y)\,\mathrm{d}S(y) \longrightarrow 0. 
\end{align*}
于是对\eqref{eq:green2}式取极限得
\begin{align}
u(x)=\int_{\partial\Omega}\left[\varGamma (y-x)\frac{\partial u}{\partial n}(y)-u(y)\frac{\partial}{\partial n}\varGamma (y-x)\right]\mathrm{d}S(y)+\int_\Omega \varGamma (y-x)f(y)\,\mathrm{d}y. \label{repr}
\end{align}

另一方面，设 $\phi^x$ 是调和函数且边界值等于 $\varGamma (y-x)$. 在 $\Omega$ 上对 $u$ 和 $\phi^x$ 应用\hyperref[Basis of Analytics-theorem:Green第二恒等式]{Green第二恒等式}，得
\begin{align}
0=\int_{\partial\Omega}\left[u(y)\frac{\partial\phi^x}{\partial n}(y)-\phi^x(y)\frac{\partial u}{\partial n}(y)\right]\mathrm{d}S(y)+\int_\Omega \phi^x(y) f(y)\,\mathrm{d}y, \label{phi}
\end{align}
其中利用了 $-\Delta u=f$ 和 $-\Delta\phi^x=0$.
将 \eqref{repr}式与 \eqref{phi} 式相加，并注意到 $\phi^x|_{\partial\Omega}=\varGamma (y-x)$，得到
\begin{align*}
u(x)
&= -\int_{\partial\Omega} u(y)\frac{\partial}{\partial n}\big(\varGamma (y-x)-\phi^x(y)\big)\,\mathrm{d}S(y)
+\int_\Omega \big(\varGamma (y-x)-\phi^x(y)\big) f(y)\,\mathrm{d}y \\
&= -\int_{\partial\Omega} g(y)\frac{\partial}{\partial n}G(x,y)\,\mathrm{d}S(y)
+\int_\Omega G(x,y) f(y)\,\mathrm{d}y,
\end{align*}
其中 $G(x,y)=\varGamma (y-x)-\phi^x(y)$. 这就是 \eqref{eq:green-sol}式.
\end{proof}

\begin{theorem}[Green 函数的对称性]
对所有 $x,y\in\Omega$，$x\neq y$，有
\begin{align*}
G(y,x) = G(x,y),
\end{align*}
其中$G(x,y)$是$\Omega$上的Green函数.
\end{theorem}

\begin{proof}
固定 $x,y\in\Omega$，$x\neq y$。取充分小 $\varepsilon>0$，使得 $B(x,\varepsilon)\cup B(y,\varepsilon)\subset\Omega$ 和 $B(x,\varepsilon)\cap B(y,\varepsilon)=\varnothing$。令
\begin{align*}
v(z) = G(x,z), \qquad w(z) = G(y,z),
\end{align*}
则
\begin{align*}
\Delta v(z) = 0,\quad (z\neq x); \qquad \Delta w(z) = 0,\quad (z\neq y).
\end{align*}
在区域 $\Omega\setminus [B(x,\varepsilon)\cup B(y,\varepsilon)]$ 上对 $v(z)$ 和 $w(z)$ 应用\hyperref[Basis of Analytics-theorem:Green第一恒等式]{Green第一恒等式} 的公式得
\begin{align}
\int_{\partial B(x,\varepsilon)} \left( \frac{\partial v}{\partial n}w - \frac{\partial w}{\partial n}v \right) \mathrm{d}S(z) = \int_{\partial B(y,\varepsilon)} \left( \frac{\partial w}{\partial n}v - \frac{\partial v}{\partial n}w \right) \mathrm{d}S(z), \label{eq:sym-identity}
\end{align}
其中 $n$ 表示 $B(x,\varepsilon)\cup B(y,\varepsilon)$ 上单位内法向量。

我们首先考虑上式 \eqref{eq:sym-identity} 的左端项。由\refthe{theorem:Dirichlet问题的解的最大模估计}知$\phi^x(z)$ 在 $\Omega$ 是有界的调和函数且 $w$ 在 $x$ 附近光滑，故
\begin{align*}
\left| \int_{\partial B(x,\varepsilon)} \frac{\partial w}{\partial n}v\,\mathrm{d}S(z) \right| \leqslant C\varepsilon^{n-1}\sup_{\partial B(x,\varepsilon)} |v|,
\end{align*}
其中 $C$ 是不依赖于 $\varepsilon$ 的常数。因此
\begin{align*}
\lim_{\varepsilon\to 0} \int_{\partial B(x,\varepsilon)} \frac{\partial v}{\partial n}v\,\mathrm{d}S(z) = 0.
\end{align*}
另外 $v(z)=\varGamma (z-x)-\phi^x(z)$，从而
\begin{align*}
\lim_{\varepsilon\to 0} \int_{\partial B(x,\varepsilon)} \frac{\partial v(z)}{\partial n}w(z)\,\mathrm{d}S(z)
= \lim_{\varepsilon\to 0} \int_{\partial B(x,\varepsilon)} \frac{\partial}{\partial n}\varGamma (z-x)w(z)\,\mathrm{d}S(z) 
= \lim_{\varepsilon\to 0} \int_{\partial B(x,\varepsilon)} w(z)\,\mathrm{d}S(z) 
= w(x).
\end{align*}
因此，当 $\varepsilon\to 0$ 时， \eqref{eq:sym-identity}式 的左端趋于 $w(x)$，而其右端趋于 $v(y)$。于是
\begin{align*}
w(x) = v(y),
\end{align*}
从而
\begin{align*}
G(y,x) = G(x,y).
\end{align*}
至此我们完成定理的证明.
\end{proof}

\subsubsection{半空间上的 Green 函数}

\begin{definition}
称函数
\begin{align*}
K(x,y)\triangleq \frac{2x_n}{n\alpha(n)|y-x|^n}\quad (x\in\mathbb{R}_+^n,\; y\in\partial\mathbb{R}_+^n=\mathbb{R}^{n-1})
\end{align*}
为上半空间 $\mathbb{R}_+^n$ 的 \textbf{Poisson 核}.
\end{definition}

\begin{definition}
设$g$ 是 $\mathbb{R}^{n-1}$ $(n\geqslant 2)$ 上有界连续函数. 对$\forall x\in\mathbb{R}_+^n$，称
\begin{align*}
u(x)\triangleq \int_{\mathbb{R}^{n-1}} K(x,y)g(y)\,\mathrm{d}y
=\frac{2x_n}{n\alpha(n)}\int_{\mathbb{R}^{n-1}}\frac{g(y)}{|y-x|^n}\,\mathrm{d}y 
\end{align*}
为上半空间 $\mathbb{R}_+^n$ 上以 $g$ 为边界值的 \textbf{Poisson 公式}或\textbf{Poisson 积分}.
\end{definition}

\begin{theorem}
假设 $g$ 是 $\mathbb{R}^{n-1}$ $(n\geqslant 2)$ 上有界连续函数，$u(x)$是上半空间 $\mathbb{R}_+^n$ 上以 $g$ 为边界值的Poisson积分,
则
\begin{enumerate}[(1)]
\item $u(x)$ 是 $\mathbb{R}_+^n$ 上无穷次可微的有界函数；

\item $\Delta u(x)=0,\ \forall x\in\mathbb{R}_+^n$；

\item 对于$\forall x_0\in\partial\mathbb{R}_+^n$，当 $x\in\mathbb{R}_+^n$ 且 $x\to x_0$ 时，$u(x)\to g(x_0)$。
\end{enumerate}
\end{theorem}
\begin{proof}
\begin{enumerate}[(1)]
\item 对$\forall x\in\mathbb{R}_+^n$，我们可以证明
\begin{align}
\int_{\mathbb{R}^{n-1}} K(x,y)\,\mathrm{d}y = 1. \label{eq:poisson-half-normal}
\end{align}
由定理的假设，存在 $M>0$，使得 $|g|\leqslant M$。因此$u(x)$ 是有界的。当 $x\neq y$ 时，$K(x,y)$ 是光滑的。我们易证 $u\in C^\infty(\mathbb{R}_+^n)$。

\item 当 $x\neq y$ 时，$\Delta_x G(x,y)=0$，于是 $\Delta_x \frac{\partial}{\partial y_n}G(x,y)=0$。所以，当 $x\in\mathbb{R}_+^n$，$y\in\partial\mathbb{R}_+^n$ 时，有 $\Delta_x K(x,y)=0$，且
\begin{align*}
\Delta u(x) = \int_{\mathbb{R}^{n-1}} \Delta_x K(x,y)g(y)\,\mathrm{d}y = 0.
\end{align*}

\item 固定 $x_0\in\partial\mathbb{R}_+^n=\mathbb{R}^{n-1}$，$\varepsilon>0$。取 $\delta_1>0$ 足够小，使得当 $y\in\mathbb{R}^{n-1}$ 且 $|y-x_0|\leqslant \delta_1$ 时，有
\begin{align*}
|g(y)-g(x_0)| \leqslant \frac{\varepsilon}{2}.
\end{align*}
于是当 $x\in\mathbb{R}_+^n$ 且 $|x-x_0|\leqslant \frac{\delta_1}{2}$ 时，利用等式 \eqref{eq:poisson-half-normal} 我们计算得
\begin{align*}
|u(x)-g(x_0)| &= \left| \int_{\mathbb{R}^{n-1}} K(x,y)[g(y)-g(x_0)]\,\mathrm{d}y \right| \\
&\leqslant \int_{\mathbb{R}^{n-1}\cap B(x_0,\delta_1)} K(x,y)|g(y)-g(x_0)|\,\mathrm{d}y \\
&\quad + \int_{\mathbb{R}^{n-1}\setminus B(x_0,\delta_1)} K(x,y)|g(y)-g(x_0)|\,\mathrm{d}y \\
&\leqslant \frac{\varepsilon}{2} + 2M\int_{\mathbb{R}^{n-1}\setminus B(x_0,\delta_1)} K(x,y)\,\mathrm{d}y,
\end{align*}
这里 $B(x_0,\delta_1)$ 是 $\mathbb{R}^{n-1}$ 上以 $x_0$ 为心，$\delta_1$ 为半径的 $n-1$ 维球。

如果 $|x-x_0|\leqslant \frac{\delta_1}{2}$ 和 $|y-x_0|\geqslant \delta_1$，我们有
\begin{align*}
|x-x_0| \leqslant \frac{1}{2}|y-x_0|,
\end{align*}
从而
\begin{align*}
|y-x_0| \leqslant |y-x| + |x-x_0| \leqslant |y-x| + \frac{1}{2}|y-x_0|,
\end{align*}
因此
\begin{align*}
|y-x| \geqslant \frac{1}{2}|y-x_0|.
\end{align*}
当 $|x-x_0|\leqslant \frac{\delta_1}{2}$ 时，我们得到
\begin{align*}
\int_{\mathbb{R}^{n-1}\setminus B(x_0,\delta_1)} K(x,y)\,\mathrm{d}y
\leqslant x_n \frac{2^{n+1}}{n\alpha(n)} \int_{\mathbb{R}^{n-1}\setminus B(x_0,\delta_1)} |y-x_0|^{-n}\,\mathrm{d}y 
\leqslant C(n,\delta_1)x_n,
\end{align*}
这里 $C(n,\delta_1)$ 是仅依赖于 $n,\delta_1$ 的常数。取 $\delta_2>0$ 足够小，使当 $0<x_n\leqslant \delta_2$ 时，有
\begin{align*}
2M\int_{\mathbb{R}^{n-1}\setminus B(x_0,\delta_1)} K(x,y)\,\mathrm{d}y \leqslant \frac{\varepsilon}{2}.
\end{align*}
因此，当 $|x-x_0|\leqslant \delta = \min\{\frac{\delta_1}{2},\delta_2\}$ 时，有
\begin{align*}
|u(x)-g(x_0)| \leqslant \varepsilon.
\end{align*}
\end{enumerate}
\end{proof}

\subsubsection{球上的 Green 函数}

\begin{definition}
设$R\in \mathbb{R}$,称函数
\begin{align*}
K(x,y)\triangleq \frac{R^2-|x|^2}{n\alpha(n)R|x-y|^n}\quad (x\in B(0,R),\; y\in\partial B(0,R))
\end{align*}
为球$B(0,R)$的\textbf{Poisson核}.
\end{definition}

\begin{definition}
设$R\in \mathbb{R}$,$g\in C(\partial B(0,R))$. 对任意 $x\in B(0,R)$，称
\begin{align}\label{eq:poisson-ball}
u(x)\triangleq \int_{\partial B(0,R)} K(x,y)g(y)\,\mathrm{d}S(y)
=\frac{R^2-|x|^2}{n\alpha(n)R}\int_{\partial B(0,R)}\frac{g(y)}{|x-y|^n}\,\mathrm{d}S(y). 
\end{align}
为球$B(0,R)$上以 $g$ 为边界值的 \textbf{Poisson公式}或\textbf{Poisson积分}.
\end{definition}

\begin{theorem}
假设 $B(0,R)\subset\mathbb{R}^n$，$g:\partial B(0,R)\to\mathbb{R}$ 连续，$u(x)$是球$B(0,R)$上以 $g$ 为边界值的Poisson积分,则
\begin{enumerate}[(1)]
\item $u(x)$ 是 $B(0,R)$ 上无穷次可微的有界函数；
\item $\Delta u(x)=0,\ x\in B(0,R)$；
\item 任给 $x_0\in\partial B(0,R)$，当 $x\in B(0,R)$ 且 $x\to x_0$ 时，$u(x)\to g(x_0)$。
\end{enumerate}
\end{theorem}
\begin{remark}
当 $n=2$ 时，在公式 \eqref{eq:poisson-ball} 中用极坐标表示 $x=(\rho\cos\theta,\rho\sin\theta)$，$y=(R\cos\varphi,R\sin\varphi)$，记 $u(x)=u(\rho,\theta)$，则公式 \eqref{eq:poisson-ball} 可以表示成
\begin{align*}
u(\rho,\theta) = \frac{1}{2\pi}\int_0^{2\pi} \frac{R^2-\rho^2}{R^2+\rho^2-2R\rho\cos(\varphi-\theta)}g(\varphi)\,\mathrm{d}\varphi,
\end{align*}
其中 $g(\varphi)=g(R\cos\varphi,R\sin\varphi)$。实际上，我们在“复变函数”这门课程中学过以上公式。
\end{remark}
\begin{proof}

对任意 $x\in B(0,R)$，我们对特殊情形 $u=1$ 应用\refthe{theorem:Dirichlet问题解的表达式}得到
\begin{align*}
\int_{\partial B(0,R)} K(x,y)\,\mathrm{d}S(y) = 1.
\end{align*}
由定理的假设，存在 $M>0$，使得 $|g|\leqslant M$。因此$u(x)$ 是有界的。当 $x\neq y$ 时，$K(x,y)$ 是光滑的，我们易证 $u\in C^\infty(B(0,R))$。

当 $x\neq y$ 时，$\Delta_x G(x,y)=0$，于是 $\Delta_x \frac{\partial}{\partial n}G(x,y)=0$。从而，当 $x\in B(0,R)$，$y\in\partial B(0,R)$ 时，$\Delta_x K(x,y)=0$ 且
\begin{align*}
\Delta u(x) = \int_{\partial B(0,R)} \Delta_x K(x,y)g(y)\,\mathrm{d}S(y) = 0.
\end{align*}

(3) 固定 $x_0\in\partial B(0,R)$，$\varepsilon>0$。取 $\delta_1>0$ 足够小，使得当 $y\in\partial B(0,R)$ 且 $|y-x_0|\leqslant \delta_1$ 时，
\begin{align*}
|g(y)-g(x_0)| \leqslant \frac{\varepsilon}{2},
\end{align*}
于是当 $x\in B(0,R)$ 且 $|x-x_0|\leqslant \frac{\delta_1}{2}$ 时，利用归一化等式我们计算得
\begin{align*}
|u(x)-g(x_0)| &= \left| \int_{\partial B(0,R)} K(x,y)[g(y)-g(x_0)]\,\mathrm{d}S(y) \right| \\
&\leqslant \int_{\partial B(0,R)\cap\{|y-x_0|\leqslant \delta_1\}} K(x,y)|g(y)-g(x_0)|\,\mathrm{d}S(y) \\
&\quad + \int_{\partial B(0,R)\cap\{|y-x_0|>\delta_1\}} K(x,y)|g(y)-g(x_0)|\,\mathrm{d}S(y) \\
&\leqslant \frac{\varepsilon}{2} + \frac{2M(R^2-|x|^2)R^{n-2}}{(\delta_1/2)^n}.
\end{align*}
取 $\delta>0$ 足够小，当 $|x-x_0|<\delta$ 时，$R^2-|x|^2$ 足够小，于是我们得到
\begin{align*}
|u(x)-g(x_0)| \leqslant \varepsilon.
\end{align*}
至此我们完成定理的证明.
\end{proof}







\end{document}